\section{Conclusion}

NEO-Watch delivers the six required features and the three optional ones we chose. A user can
create an account, see what will pass close to the Earth in the coming days, search the objects
and export them, follow favourites over time and be warned when an object matches their own
criteria; an administrator manages the accounts and the data, which is kept up to date
automatically.

\vspace{0.25cm}

The main technical difficulty was not where we expected it. The features themselves were quick
to write once the data model was fixed; the hard part was the integration: making a long NASA
import live with an API that must keep answering, sharing a database between threads, handling
the rate limit of the free NASA key, and making the backend and the frontend agree on every
answer. Testing the application end to end, and not only each layer alone, is what revealed most
of the bugs.

\subsection{Limits}

\begin{itemize}
  \item NEO-Watch only displays and organises NASA's data: it computes no orbit and no impact
  probability, and its content is only as fresh and as accurate as the NeoWs API.
  \item With the free \texttt{DEMO\_KEY}, the first import takes several hours, since it is cut
  by the rate limit; a personal NASA key is needed for a real use.
  \item NEOs declared by users are not moderated, unlike what the analysis report planned.
  \item The notifications are only shown inside the application, when the user logs in.
  \item Passwords are hashed with SHA-256: a slow algorithm made for passwords (bcrypt, Argon2)
  would resist brute force better.
\end{itemize}

\subsection{Possible improvements}

\begin{itemize}
  \item Moderation of the declared NEOs through the \texttt{proposal} table, which already exists.
  \item Notifications by e-mail, sent after each daily import.
  \item A user page to change one's e-mail address and password.
  \item Tests of the Streamlit pages, which are only tested by hand today.
\end{itemize}
